\section{Algorithms, complexity, and numerical considerations}
\label{sec:algorithms}

Section~\ref{sec:dp} derives the exact dynamic program (DP) that sits at the core of BayesBreak.
The present section complements that development with implementation-ready pseudocode and a
collection of numerical-stability considerations.
The primary purpose is pragmatic: to make explicit what must be computed, stored, and
normalized in order to reproduce the theoretical quantities (evidences, boundary posteriors,
Bayes curves) without relying on undocumented conventions.

\subsection{Block-evidence precomputation via prefix sums}
\label{sec:alg-blocks}

For conjugate exponential-family blocks, Theorem~\ref{thm:ef-integral} shows that the single
block evidence and moments can be expressed as functions of aggregated sufficient statistics.
This enables an \emph{exact} $\mathcal{O}(n^2)$ precomputation of all candidate blocks $(i,j]$
with $0\le i<j\le n$. Throughout this section we use the body convention: a block $(i,j]$ has
length $j-i$ and contains the indices $i+1,\dots,j$.

In practice, it is numerically preferable to store \emph{log} block evidences,
\begin{equation}
    \ell_{i j}^{(0)} \;\coloneqq\; \log \widetilde{A}^{(0)}_{i j},
\end{equation}
and similarly for positive moment numerators required for the Bayes regression curve. If a chosen
moment numerator can be zero or negative, implementations use tagged signed-log storage rather than
silently taking a logarithm: a signed quantity $m_{ij}$ is represented as $(s_{ij},a_{ij})$ with
$s_{ij}\in\{-1,0,1\}$ and $a_{ij}=\log |m_{ij}|$ when $s_{ij}\ne0$, while $s_{ij}=0$ records exact
zero. The output contract for each \texttt{BlockRoutine} is explicit: return a finite log evidence
for admissible blocks, $-\infty$ for excluded blocks, optional positive log-moment numerators, and
optional signed-log moment records with this zero convention. Downstream DP code must reject
untagged nonpositive moment numerators rather than treating them as log values. Algorithm~\ref{alg:block-precompute} expresses the precomputation at an abstract level, with a
single call to a \texttt{BlockRoutine} that implements the family-specific closed forms given in
Section~\ref{sec:families}.

\begin{algorithm}[H]
\caption{PrecomputeBlocks$(y,x,w,\texttt{BlockRoutine})$}
\label{alg:block-precompute}
\KwIn{Ordered observations $\{(x_i,y_i,w_i)\}_{i=1}^n$ (with $w_i\equiv 1$ if unweighted);
family-specific routine \texttt{BlockRoutine} implementing Theorem~\ref{thm:ef-integral}.}
\KwOut{Matrices of log-block evidences $\ell^{(0)}_{ij}=\log \widetilde{A}^{(0)}_{ij}$ on the index set $\{0\le i<j\le n\}$ and
(optional) moment numerators, stored as $\log \widetilde{A}^{(r)}_{ij}$ only when they are positive and otherwise in a signed-log or linear representation.}
\BlankLine
\textbf{1. Prefix summaries}\;
Compute prefix cumulative sufficient statistics and weights required by \texttt{BlockRoutine}:
store $W_j\coloneqq\sum_{t=1}^j w_t$, $S_j\coloneqq\sum_{t=1}^j w_t\,T(y_t)$, and any additional
family-specific terms (e.g., $Q_j\coloneqq\sum_{t=1}^j w_t y_t^2$ for Gaussian). Set $W_0=S_0=Q_0=0$.
\BlankLine
\textbf{2. Enumerate blocks}\;
\For{$i\leftarrow 0$ \KwTo $n-1$}{
  \For{$j\leftarrow i+1$ \KwTo $n$}{
    Extract block summaries for $y_{i+1:j}$ via $W_{ij}=W_j-W_i$, $S_{ij}=S_j-S_i$, $Q_{ij}=Q_j-Q_i$ (constant time);
    \texttt{BlockRoutine} returns $\log A^{(0)}_{ij}$ and any moment numerators $\log A^{(r)}_{ij}$;
    Add any length factor $\log g(\Delta_x(i,j))$ and hazard constants as in \eqref{eq:Atilde};
    Store $\ell^{(0)}_{ij}\leftarrow \log \widetilde{A}^{(0)}_{ij}$ and, if needed,
    $\ell^{(r)}_{ij}\leftarrow \log \widetilde{A}^{(r)}_{ij}$.
  }
}
\end{algorithm}

\paragraph{Memory notes.}
Storing all $\ell^{(0)}_{ij}$ uses $\Theta(n^2)$ memory.
When $n$ is large but $k_{\max}$ is moderate, the memory footprint is reduced by streaming
blocks in the inner DP loop (computing $\ell^{(0)}_{ij}$ on the fly from prefix statistics), at
cost of recomputing some block summaries.
The DP theory remains unchanged.

As sanity checks, verify that invalid blocks are represented by $-\infty$ log-scores, that $\log C_k$ is computed with the same admissibility rules as the DP, that forward and backward full-sequence evidences agree, and that boundary-event marginals sum to $k-1$ conditional on fixed $k$.

\subsection{Log-space dynamic programming}
\label{sec:alg-dp-impl}

The forward/backward DP in Section~\ref{sec:dp} sums products of block evidences.
To avoid underflow/overflow, work in log space and use a stable
\texttt{logsumexp} primitive.
Algorithm~\ref{alg:dp-log} presents a log-space version of the recursions in
\eqref{eq:LR}--\eqref{eq:post-k}.

\begin{algorithm}[H]
\caption{DPLog$(\ell^{(0)}_{ij},k_{\max},\log p(k),\log C_k)$}
\label{alg:dp-log}
\KwIn{Log block evidences $\ell^{(0)}_{ij}=\log \widetilde{A}^{(0)}_{ij}$ on the index set $\{0\le i<j\le n\}$ (block $(i,j]$);
maximum segments $k_{\max}$;
log prior on $k$, $\log p(k)$;
log normalizers $\log C_k$ computed under the same admissible-block and length-prior convention as the score array. For the index-uniform fixed-$k$ prior with $g\equiv1$, $C_k=\binom{n-1}{k-1}$ unless the caller has intentionally absorbed a normalized prior into the block or count scores.}
\KwOut{Forward messages $\log \widetilde{L}_{k,j}$, backward messages $\log \widetilde{R}_{k,i}$,
log evidence per $k$, and normalized posterior $p(k\mid y)$.}
\BlankLine
\textbf{1. Initialize}\;
Set $\log \widetilde{L}_{0,0}\leftarrow 0$ and $\log \widetilde{L}_{0,j}\leftarrow -\infty$ for $j\ge 1$;
set $\log \widetilde{R}_{0,n}\leftarrow 0$ and $\log \widetilde{R}_{0,i}\leftarrow -\infty$ for $i\le n-1$.
\BlankLine
\textbf{2. Forward recursion}\;
\For{$k\leftarrow 1$ \KwTo $k_{\max}$}{
  \For{$j\leftarrow k$ \KwTo $n$}{
    $\log \widetilde{L}_{k,j}\leftarrow \texttt{logsumexp}_{h\in\{k-1,\dots,j-1\}}\big(\log \widetilde{L}_{k-1,h} + \ell^{(0)}_{h,j}\big)$;
  }
}
\BlankLine
\textbf{3. Backward recursion}\;
\For{$k\leftarrow 1$ \KwTo $k_{\max}$}{
  \For{$i\leftarrow n-k$ \KwTo $0$ (decreasing)}{
    $\log \widetilde{R}_{k,i}\leftarrow \texttt{logsumexp}_{h\in\{i+1,\dots,n-k+1\}}\big(\ell^{(0)}_{i,h} + \log \widetilde{R}_{k-1,h}\big)$;
  }
}
\BlankLine
\textbf{4. Evidence and posterior over $k$}\;
For each $k\in\{1,\dots,k_{\max}\}$ set
$\log p(y\mid k)\leftarrow \log \widetilde{L}_{k,n} - \log C_k$;
compute unnormalized log posterior scores
$\log \pi_k \leftarrow \log p(k) + \log p(y\mid k)$;
set $p(k\mid y)\leftarrow \exp\{\log \pi_k - \texttt{logsumexp}_{k'}(\log \pi_{k'})\}$.
\end{algorithm}

The backward upper endpoint $n-k+1$ is the last split that leaves $k-1$ nonempty suffix blocks to
cover indices $h+1,\dots,n$. Empty summation ranges are interpreted as $-\infty$ in log space, so
edge cases are handled by the same admissibility convention as the forward recursion.

\paragraph{Boundary marginals and Bayes curves.}
Once $\log\widetilde{L}$ and $\log\widetilde{R}$ are available, the boundary marginal in
\eqref{eq:boundary-post} can be computed stably by subtracting the common normalizer
$\log\widetilde{L}_{k,n}$.
Similarly, Bayes regression moments (Theorem~\ref{thm:dp-correctness}, Step~6; cf.\ \eqref{eq:segmom})
can be computed by replacing $\ell^{(0)}_{ij}$ with compatible moment numerators and reusing the
same log-space DP bookkeeping. Signed or nonpositive moment targets require the signed-log or
scaled-linear convention noted in Section~\ref{sec:alg-blocks}.
A signed-log accumulator stores each target as a sign bit and a log-magnitude; sums of signed
terms are evaluated by separately applying \texttt{logsumexp} to positive and negative parts
and subtracting the two scaled results only at the final reporting step. This convention is part
of the output contract for any routine that returns moment numerators rather than evidences.
The backward-loop bound $h\le n-k+1$ in Algorithm~\ref{alg:dp-log} enforces the nonempty-suffix
convention: after choosing the next endpoint $h$ for the block $(i,h]$, the remaining $k-1$
blocks must occupy at least one observation each in $(h,n]$. Thus the summation is empty, and the
message is $-\infty$, whenever the remaining suffix cannot support the requested number of blocks.

\subsection{MAP segmentation and backtracking}
\label{sec:alg-map}

The DP used for evidences can be augmented to recover a MAP segmentation in
$\mathcal{O}(k_{\max}n^2)$ time.
During the forward recursion, one stores, for each $(k,j)$, the maximizing split point
\begin{equation}
    h^\star(k,j)\in\arg\max_{h\in\{k-1,\dots,j-1\}} \Big\{M_{k-1,h}+\ell^{(0)}_{h,j}\Big\},
    \label{eq:map-backpointer}
\end{equation}
where $M_{k,j}$ is the max-sum message defined in \eqref{eq:max-sum-inline}.
Algorithm~\ref{alg:map-forward} computes $M_{k,j}$ and the backpointers; Algorithm~\ref{alg:map-backtrack} then backtracks from $(\widehat{k},n)$ to recover the boundary vector
$\widehat{b}=(\widehat{b}_0,\dots,\widehat{b}_{\widehat{k}})$.

\begin{theorem}[Correctness of max-sum DP for joint MAP segmentation]
\label{thm:map-correctness}
For any choice of segment count $\widehat k\in\{1,\dots,k_{\max}\}$, Algorithms~\ref{alg:map-forward}
and~\ref{alg:map-backtrack} jointly return the (or a) boundary vector
$\widehat b=(\widehat b_0,\dots,\widehat b_{\widehat k})$ that attains
\[
\widehat b \in \arg\max_{t\in T_{\widehat k,n}}\Big\{
   \log p(\widehat k) - \log C_{\widehat k} + \sum_{q=1}^{\widehat k} \ell^{(0)}_{t_{q-1},t_q}
\Big\},
\]
where $T_{\widehat k,n}$ is the set of admissible boundary vectors of length $\widehat k$ on $n$
observations under the segment-factorized partition prior of Section~\ref{sec:dp}. In particular,
$\widehat b$ is a joint MAP segmentation conditional on the chosen $\widehat k$, and any tie among
maximizers is broken deterministically by the back-pointer convention encoded in
$h^\star(k,j)$. The runtime is $\mathcal{O}(k_{\max}n^2)$ and the working space is
$\mathcal{O}(k_{\max}n)$.
\end{theorem}

A proof is in Appendix~\ref{app:max-sum-proof}; it proceeds by induction on $k$ using the
sub-segmentation optimality property of the max-sum recursion.

\begin{algorithm}[H]
\caption{MAPForward$(\ell^{(0)}_{ij},k_{\max})$}
\label{alg:map-forward}
\KwIn{Log block evidences $\ell^{(0)}_{ij}$ on $\{0\le i<j\le n\}$; $k_{\max}$.}
\KwOut{Max-sum messages $M_{k,j}$ and backpointers $h^\star(k,j)$ for $k\le k_{\max}$.}
\BlankLine
Set $M_{0,0}\leftarrow 0$, $M_{0,j}\leftarrow -\infty$ for $j\ge 1$\;
\For{$k\leftarrow 1$ \KwTo $k_{\max}$}{
  \For{$j\leftarrow k$ \KwTo $n$}{
    $h^\star(k,j)\leftarrow \arg\max_{h\in\{k-1,\dots,j-1\}}\{M_{k-1,h}+\ell^{(0)}_{h,j}\}$\;
    $M_{k,j}\leftarrow M_{k-1,h^\star(k,j)}+\ell^{(0)}_{h^\star(k,j),j}$\;
  }
}
\end{algorithm}

\begin{algorithm}[H]
\caption{MAPBacktrack$(h^\star,\widehat{k})$}
\label{alg:map-backtrack}
\KwIn{Backpointers $h^\star(k,j)$ as in \eqref{eq:map-backpointer}; chosen segment count $\widehat{k}$.}
\KwOut{MAP boundary vector $\widehat{b}_0=0<\widehat{b}_1<\cdots<\widehat{b}_{\widehat{k}}=n$.}
\BlankLine
Set $\widehat{b}_{\widehat{k}}\leftarrow n$ and $q\leftarrow\widehat{k}$.
\While{$q\ge 1$}{
  $\widehat{b}_{q-1}\leftarrow h^\star(q,\widehat{b}_{q})$\;
  $q\leftarrow q-1$\;
}
Return $\widehat{b}$.
\end{algorithm}

\subsection{Pooling across sequences and groups}
\label{sec:alg-pooling}

For replicated or multi-subject data with shared boundaries (Section~\ref{sec:replicates}),
Theorem~\ref{thm:multisubject} implies that subject-level block evidences multiply.
In log space, pooling is therefore a simple sum:
\begin{equation}
    \ell^{(0,\text{pooled})}_{ij} \;=\; \sum_{s=1}^S \ell^{(0,s)}_{ij}.
\end{equation}
The same applies to known groups, with the sum taken over subjects in a group.
The resulting pooled matrices can be passed unchanged to Algorithm~\ref{alg:dp-log}.

\begin{algorithm}[H]
\caption{Pooled-DP for shared-boundary multi-subject inference (log-space)}
\label{alg:pool-shared-boundary}
\KwIn{Per-subject log-block-evidence arrays $\{\ell^{(0,s)}_{ij}\}_{s=1}^S$ for $0\le i<j\le n$;
optional length-prior log-cohesion $\log g(\Delta_x(i,j))$; $k_{\max}$; group-label set
$\{z_s\}_{s=1}^S$ if pooling by known group, else $z_s\equiv 1$ for global pooling.}
\KwOut{Per-group pooled posteriors $P_g(k\mid y)$, $P_g(t_p\mid y,k)$, MAP segmentations
$\widehat t^{\mathrm{MAP}}_g(\widehat k_g)$, and segment-conditional parameter posteriors.}
\For{each distinct group label $g\in\{z_s\}$}{
  Let $\mathcal{S}_g=\{s:z_s=g\}$\;
  $\ell^{(0,\mathrm{pool})}_{ij,g}\leftarrow \sum_{s\in\mathcal{S}_g}\ell^{(0,s)}_{ij}$ for every admissible $(i,j]$\;
  \tcp{Length prior absorbed once at the pooled level.}
  $\widetilde\ell^{(0,\mathrm{pool})}_{ij,g}\leftarrow \ell^{(0,\mathrm{pool})}_{ij,g}+\log g(\Delta_x(i,j))$\;
  Run Algorithm~\ref{alg:dp-log} with score $\widetilde\ell^{(0,\mathrm{pool})}_{ij,g}$ to obtain $P_g(k\mid y)$, $P_g(t_p\mid y,k)$, $\widehat k_g$\;
  Run the max-sum routine (Algorithms~\ref{alg:map-forward}--\ref{alg:map-backtrack}) on the same score to obtain $\widehat t^{\mathrm{MAP}}_g(\widehat k_g)$\;
  Conditional on $\widehat t^{\mathrm{MAP}}_g$, recover subject-specific segment-parameter posteriors from each $s\in\mathcal{S}_g$ block update\;
}
\Return $\{P_g(k\mid y),P_g(t_p\mid y,k),\widehat t^{\mathrm{MAP}}_g,\text{per-subject parameter posteriors}\}_g$\;
\tcp{Complexity: $\Theta\!\bigl(\sum_g(|\mathcal{S}_g|n^2 + k_{\max}n^2)\bigr)$ time and
$\Theta(n^2 + k_{\max}n)$ working memory per group when subject sums are streamed.}
\tcp{Numerical: pool in log space and use \texttt{logsumexp}; carry per-subject log-blocks to
floating-point underflow safely; sum $-\infty$ contributions are absorbed at the pooled step.}
\end{algorithm}

\subsection{Latent-group EM implementation notes}
\label{sec:alg-em}

Section~\ref{sec:latent-em} formulates latent groups as a \emph{template mixture}: each group has one
exported segmentation template, while subject-specific segment parameters have already been integrated out in
subject-level block evidences. From an implementation perspective, the M-step for each group is therefore a
standard max-sum BayesBreak backtracking run with a responsibility-weighted score matrix.
If $r_{sg}$ denotes the responsibility of subject $s$ for group $g$, let
$n_g=\sum_s r_{sg}$. The effective block score for the penalized template objective is
\begin{equation}
    B^{(g)}_{ij}
    \;=\; n_g\log g(\Delta_x(i,j)) + \sum_{s=1}^S r_{sg}\, \ell^{(0,s)}_{ij},
\end{equation}
where $\ell^{(0,s)}_{ij}=\log A^{(0,s)}_{ij}$ is the subject-level log-block evidence.
Running the max-sum DP over this score matrix for each $k$ and then adding the responsibility-scaled
count offset $n_g\{\log p(k)-\log C_k\}$ yields the exact template update for group $g$ under the
objective in Section~\ref{sec:latent-em}.
The monotonicity guarantee in Theorem~\ref{thm:em-monotone} applies precisely because these E- and M-step
updates are exact for the stated template-mixture objective.

\begin{algorithm}[H]
\caption{Latent-template EM with exact responsibility and exact max-sum M-step (log-space)}
\label{alg:latent-em-detail}
\KwIn{Per-subject log-block-evidence arrays $\{\ell^{(0,s)}_{ij}\}_{s=1}^S$; group count $G$;
$k_{\max}$; restart count $R\ge 5$; mixing-weight floor $\pi_{\min}\in[0,1)$ (optional);
log-cohesion $\log g(\Delta_x(\cdot))$ and partition-prior $\log p(k)/C_k$.}
\KwOut{Mixing weights $\widehat\pi\in\Delta^{G-1}$; templates $\{\widehat\tau_g\}_{g=1}^G$; final
responsibilities $\{r_{sg}\}$; achieved finite-mixture objective $\ell_\star$.}
\For{restart $r=1$ \KwTo $R$}{
  Initialize $(\pi^{(0)},\tau^{(0)})$ by (\emph{r=1}) $k$-means on pooled per-subject scores, (\emph{r=2,\dots,R}) draws from the partition prior\;
  \Repeat{$|\Delta\ell_\star|<10^{-6}$ \emph{and} templates stop changing under deterministic tie-breaking}{
    \tcp{E-step (exact responsibility update; closed-form for the current finite templates).}
    $r_{sg}\leftarrow \dfrac{\pi_g\,S_g(y^{(s)};\tau_g)}{\sum_{h=1}^G\pi_h\,S_h(y^{(s)};\tau_h)}$ for each $s,g$\;
    \tcp{M-step (i): mixing-weight update.}
    $n_g\leftarrow \sum_s r_{sg}$; $\pi_g\leftarrow \max(n_g/S,\pi_{\min})$ then renormalize\;
    \tcp{M-step (ii): exact responsibility-weighted max-sum template update.}
    \For{$g=1$ \KwTo $G$}{
      Form $B^{(g)}_{ij}=n_g\log g(\Delta_x(i,j))+\sum_s r_{sg}\ell^{(0,s)}_{ij}$\;
      Run max-sum DP + backtracking on $B^{(g)}$ with offset $n_g\{\log p(k)-\log C_k\}$ to obtain $\widehat\tau_g=(\widehat k_g,\widehat t^{(g)})$\;
    }
    Compute $\ell_\star^{(t+1)}=\sum_{s=1}^S\log\sum_g \pi_g S_g(y^{(s)};\tau_g)$\;
  }
  Record $(\pi^{(r)},\tau^{(r)},\ell_\star^{(r)})$\;
}
\Return restart with largest achieved $\ell_\star^{(r)}$\;
\tcp{Complexity (per iteration): $\Theta(GS n^2 + G k_{\max} n^2)$ time and $\Theta(Sn^2+Gk_{\max}n)$ working memory.}
\tcp{Numerical: do every responsibility update with a \texttt{logsumexp}-style normalizer; report
$\pi_{\min}>0$ as a regularized variant (Remark on regularized EM in Section~\ref{sec:latent-em}).}
\end{algorithm}

\subsection{Complexity summary and stability}
\label{sec:alg-complexity}


\begin{table}[H]
\centering
\small
\resizebox{\textwidth}{!}{%
\begin{tabular}{@{}llll@{}}
\toprule
Component & Time & Memory & Numerical note \\
\midrule
Conjugate block precomputation & $\Theta(n^2)$ & $\Theta(n^2)$ & prefix summaries; invalid blocks stored as $-\infty$ \\
Exact DP, fixed score matrix & $\Theta(k_{\max}n^2)$ & $\Theta(k_{\max}n)$--$\Theta(k_{\max}n^2)$ & logsumexp normalization \\
Boundary marginals and Bayes curve & $\Theta(k_{\max}n^2)$ & $\Theta(k_{\max}n^2)$ if all block covers are retained & common evidence normalizer \\
MAP backtracking & $\Theta(k_{\max}n^2)$ & $\Theta(k_{\max}n)$ plus backpointers & deterministic tie rule \\
Non-conjugate precomputation & $\Theta(n^2T_{\mathrm{block}})$ & $\Theta(n^2)$ plus routine state & routine convergence/failure flags \\
Shared-boundary pooling & $\Theta(Sn^2+k_{\max}n^2)$ after subject blocks & $\Theta(Sn^2)$ or streamed subject sums & sum subject log scores \\
Latent-template EM & per iteration $\Theta(GSn^2+Gk_{\max}n^2)$ & $\Theta(Sn^2+Gk_{\max}n)$ & log-responsibility normalization \\
Prediction scoring & $\Theta(Gm)$ after segment assignment; $\Theta(Gk_{\max}m^2)$ for resegmentation & problem dependent & log predictive scores \\
\bottomrule
\end{tabular}%
}
\caption{Implementation-level complexity summary. This is a theoretical accounting table; no source CSV is associated with it.}
\label{tab:complexity-summary}
\end{table}

\begin{proposition}[BayesBreak time and space complexity]
\label{prop:bb-complexity}
Suppose the chosen exponential-family block model admits constant-time block-evidence evaluation
from cumulative sufficient statistics (Sections~\ref{sec:ef}--\ref{sec:families}), and let
$n$ denote the number of observations and $k_{\max}$ the maximum segment count considered by the
DP. Then:
\begin{enumerate}
\item Block-evidence precomputation (Algorithm~\ref{alg:block-precompute}) runs in $\Theta(n^2)$ time and stores $\Theta(n^2)$ values.
\item The forward DP (Algorithm~\ref{alg:dp-log}) runs in $\Theta(k_{\max}n^2)$ time. The required working memory is $\Theta(k_{\max}n)$ if only $p(y\mid k)$ and $p(k\mid y)$ are needed, and $\Theta(k_{\max}n^2)$ if boundary marginals or backpointers must be retained.
\item The max-sum recursion plus backtracking (Algorithms~\ref{alg:map-forward}--\ref{alg:map-backtrack}) runs in $\Theta(k_{\max}n^2)$ time and $\Theta(k_{\max}n)$ working space.
\item For non-conjugate block models with per-block approximation cost $T_{\mathrm{block}}$ (e.g.\ Laplace, JJ, EP, PG mean field, 1D quadrature), block precomputation costs $\Theta(n^2 T_{\mathrm{block}})$ and the rest of the DP is unchanged.
\end{enumerate}
\end{proposition}

\begin{proof}
Parts (i)--(iii) are standard counting arguments under the constant-time
block evaluation hypothesis; the full complexity argument is given in
Appendix~\ref{app:complexity-proofs}. Part (iv) follows because the DP layer
treats block evidences as opaque inputs, so the only complexity change is in
the precomputation pass.
\end{proof}


Two numerical practices are essential in realistic data sets:
(i) always store evidences and DP messages in log space and normalize via \texttt{logsumexp};
(ii) standardize/center sufficient statistics when available (e.g., for Gaussian blocks) to
reduce catastrophic cancellation in moment computations.
These practices do not affect correctness, but they materially improve stability when $n$ is
in the hundreds or thousands.

\subsection{A numerical-implementation checklist}
\label{sec:alg-checklist}

The following checklist distills the common pitfalls we have encountered and translates them into
implementation tests. Each item should be enforced either by a unit test, an assertion in the
runtime path, or a recorded metadata check in the experiment pipeline.

\begin{enumerate}\itemsep2pt
\item \textbf{Log-space storage test.} Store all block evidences, moment numerators, and DP
messages as $\log$ values from the start, with explicit tags for signed quantities. Convert to
linear space only for display.
\item \textbf{Stable-logsumexp test.} Implement \texttt{logsumexp} as
$\mathrm{lse}(x)=m+\log\sum\exp(x_i-m)$ with $m=\max_i x_i$ \citep{blanchard2021lse}. Cache the
maximum once and reuse for both numerator and denominator when normalizing.
\item \textbf{Prefix-array test.} Compute prefix sums of $w_t$, $w_t T(y_t)$, and
(for Gaussian) $w_t y_t^2$ \emph{before} enumerating blocks. Store them as double-precision arrays
of length $n+1$ initialized at index 0. Block summaries then cost one subtraction per block.
\item \textbf{DP-invariant test.} After the forward recursion, verify that
$\log\widetilde L_{0,0}=0$ and $\log\widetilde L_{k,j}=-\infty$ for $j<k$; after the backward
recursion, verify the analogous boundary conditions on $\log\widetilde R$. Verify that
$\log\widetilde L_{k,n}=\log\widetilde R_{k,0}$ up to floating-point tolerance, that
$\sum_k p(k\mid y)=1$, that the fixed-$k$ boundary marginal sums to $k-1$, and that the reported
$C_k$ matches the prior convention used to form the block scores.
\item \textbf{Moment-scaling test.} When computing Bayes regression curves via \eqref{eq:segmom},
subtract a block-specific shift (e.g., the block posterior-mean estimate) before exponentiating
to reduce dynamic range, then add the shift back at the end.
\item \textbf{Degenerate-block test.} Guard against $W_{ij}=0$ (missing-data blocks) by returning
$\log A^{(0)}_{ij}=-\infty$ or a user-specified fallback. Guard against $Q_{ij}$ numerically
negative (from catastrophic cancellation) by clamping to zero.
\item \textbf{Reproducibility-metadata test.} Fix a seed for EM initializations and record it in the run
metadata; record hyperparameters, $k_{\max}$, the count prior $p(k)$, length factor $g$, approximation
method and settings, data-preprocessing hashes, and the exact version of every family-specific
\texttt{BlockRoutine} used.
\item \textbf{Unit tests tied to theory.} Test the block-routine output contract, the forward/backward
identity, MAP backtracking, fixed-$k$ boundary sums, signed-log moment reconstruction, and every
reported complexity mode on small instances where exhaustive enumeration is possible.
\end{enumerate}
